\subsection{Adaptive Momentum}

Newer advances in evolutionary theory, however, provide rich opportunities for new advances in evolutionary computation.
One recent advancement, the theory of adaptive momentum \citep{bohm2024reduced}, is so named because it explains the tendency for populations that have not yet converged to discover beneficial mutations at a faster rate than populations that have converged.

When a population is in the midst of a non-equilibrium transition, such as the selective sweep of a beneficial mutation, a territory range expansion, or a carrying capacity increase, some lineages experience weakened purifying selection.
If a deleterious mutation arises on one of those lineages, it is less likely to be purged.
Once the population has completed its transition, purifying selection acts with equal strength on all lineages once again.
Therefore, from the same genetic starting point, a population experiencing one of these transitions will more easily escape a local optimum than a converged population.
Furthermore, since the selective sweep of a beneficial mutation is one of these transitions, a cascade of local optima escapes is possible, leading to the appearance of adaptive momentum.
One surprising feature of adaptive momentum is that in multidimensional landscapes, adaptation along one axis can promote discoveries along others.
This last property of adaptive momentum makes it extremely potent in high-dimensional fitness landscapes.
